% Shared preamble for the Generative AI (CS 782/882, Fall 2026) lecture notes.
% Each weekNN-notes.tex file inlines this block so it compiles on its own.
\documentclass[11pt]{article}
\usepackage[margin=1in]{geometry}
\usepackage[T1]{fontenc}
\usepackage[utf8]{inputenc}
\usepackage{lmodern}
\usepackage{amsmath,amssymb,amsthm,mathtools,bm}
\usepackage{enumitem}
\usepackage{graphicx}
\usepackage{booktabs}
\usepackage{xcolor}
\usepackage{tikz}
\usetikzlibrary{arrows.meta,positioning,calc,shapes}
\usepackage{listings}
\usepackage[colorlinks=true,linkcolor=blue!60!black,urlcolor=blue!60!black,citecolor=blue!60!black]{hyperref}

\newcommand{\E}{\mathbb{E}}
\newcommand{\Var}{\mathrm{Var}}
\newcommand{\KL}{\mathrm{KL}}
\newcommand{\given}{\,\vert\,}
\newcommand{\R}{\mathbb{R}}
\DeclareMathOperator*{\argmax}{arg\,max}
\DeclareMathOperator*{\argmin}{arg\,min}
\newtheorem{theorem}{Theorem}
\newtheorem{lemma}{Lemma}
\newtheorem{definition}{Definition}
\theoremstyle{remark}
\newtheorem*{remark}{Remark}

\lstset{basicstyle=\ttfamily\small,breaklines=true,frame=single,columns=fullflexible}

% Additional definitions for this supplementary note
\newcommand{\highlight}[1]{\colorbox{yellow!40}{#1}}

\title{Supplementary Notes: DeepSeek Architecture and Training}
\author{Hong Qin, Old Dominion University}
\date{Fall 2026}

\begin{document}
\maketitle

\begin{quote}\small\itshape Note. The material in these supplementary notes was covered in the 16-week version of this course and is \textbf{optional reading} for the 8-week offering. It assumes the Week 6 lecture notes on transformers (self-attention, positional encoding, the transformer block, and next-token training) and is provided for students who want to see how those ideas are extended in current research systems.\end{quote}

\tableofcontents
\newpage

\section{Brief Introduction to the DeepSeek Model Family}

DeepSeek is a family of large language models (LLMs) designed to deliver competitive or superior performance to other frontier models while being significantly more efficient in training and inference.

Recent versions such as \emph{DeepSeek-V3} and \emph{DeepSeek-R1} are built on a decoder-only Transformer backbone but introduce three key architectural ideas:

\begin{enumerate}[label=\arabic*)]
  \item \textbf{Multi-Head Latent Attention (MLA):} instead of storing a separate key and value tensor per attention head, DeepSeek compresses keys and values into a low-dimensional \emph{latent} space and only caches these latent K/V tensors. At inference time, each head reconstructs its own keys and values from the shared latent representation. This drastically reduces KV-cache memory and bandwidth, enabling long contexts and large models with lower cost.

  \item \textbf{DeepSeekMoE (Mixture-of-Experts):} the usual dense feed-forward network (FFN) is replaced by a sparse Mixture-of-Experts: many small expert FFNs plus a routing network that selects only a few experts per token. Variants such as fine-grained experts, shared experts, and improved load balancing give very high parameter capacity while keeping the per-token compute budget modest.

  \item \textbf{Rotary Positional Embeddings (RoPE):} positional information is injected by applying position-dependent rotations to query and key vectors. In DeepSeek, RoPE is combined with MLA by splitting the representation into RoPE-aware and non-RoPE (NoPE) components, giving the model a flexible way to represent both positional and pure semantic information.

    \item \textbf{Multi-Token Generation (MTG):} DeepSeek augments standard next-token prediction with mechanisms that allow the model to propose and verify several future tokens in a single forward pass. This includes parallel multi-token prediction heads during training and speculative or Medusa-style multi-token decoding during inference. MTG reduces autoregressive latency, increases throughput, and accelerates long-form reasoning and reinforcement-learning rollouts in DeepSeek-R1.
\end{enumerate}

At a high level, a single DeepSeek Transformer block can be seen as:
\[
\text{Block}(X) = X + \text{MLA}\big(\text{LN}(X)\big)
                 + \text{MoE}\big(\text{LN}(X + \text{MLA}(\text{LN}(X)))\big),
\]
where \highlight{LN denotes Layer Normalization, MLA is the attention module with latent K/V and RoPE}, and \highlight{MoE is the sparsely activated mixture-of-experts module}.


\section{High-Level Architecture}

Let:
\begin{itemize}
  \item $L$ be the sequence length.
  \item $d_{\text{model}}$ be the model dimension.
  \item $H$ be the number of attention heads.
  \item $d_h = d_{\text{model}} / H$ be the per-head dimension.
\end{itemize}

Given a layer input
\[
X \in \mathbb{R}^{L \times d_{\text{model}}},
\]
a DeepSeek block proceeds in two main stages:
\begin{enumerate}[label=\arabic*)]
  \item \textbf{Multi-Head Latent Attention (MLA)} with RoPE:
  \[
  A = \text{MLA}( \text{LN}(X) ) \in \mathbb{R}^{L \times d_{\text{model}}}.
  \]
  \item \textbf{Mixture-of-Experts (MoE) FFN:}
  \[
  M = \text{MoE}( \text{LN}(X + A) ) \in \mathbb{R}^{L \times d_{\text{model}}}.
  \]
\end{enumerate}
The output of the block is then:
\[
Y = X + A + M.
\]

\textbf{Plain-language view:}
\begin{itemize}
  \item The attention part decides ``what each token should look at'' in the context, using MLA+RoPE efficiently.
  \item The MoE part applies several sparsely selected feed-forward ``experts'' to transform each token in a content-dependent way.
\end{itemize}

\section{Review: Standard Transformer Self-Attention}

For comparison, standard multi-head attention (MHA) proceeds as follows.

\subsection{Mathematical Formulation}

Input:
\[
X \in \mathbb{R}^{L \times d_{\text{model}}}.
\]
Linear projections:
\begin{align*}
Q &= X W_Q, \\
K &= X W_K, \\
V &= X W_V,
\end{align*}
where $W_Q, W_K, W_V \in \mathbb{R}^{d_{\text{model}} \times d_{\text{model}}}$.

We split $Q, K, V$ into $H$ heads:
\[
Q = [Q^{(1)}, \dots, Q^{(H)}], \quad
Q^{(h)} \in \mathbb{R}^{L \times d_h},
\]
and similarly for $K^{(h)}, V^{(h)}$.

For head $h$:
\[
\text{Attn}^{(h)}(X)
=
\text{softmax}\!\left(
\frac{Q^{(h)} (K^{(h)})^{\top}}{\sqrt{d_h}}
\right) V^{(h)}.
\]

We then concatenate and project:
\[
\text{MHA}(X)
= \big[ \text{Attn}^{(1)}, \dots, \text{Attn}^{(H)} \big] W_O,
\]
with $W_O \in \mathbb{R}^{d_{\text{model}} \times d_{\text{model}}}$.

\subsection{Plain-Language Explanation}

\begin{enumerate}[label=\arabic*)]
  \item For each token, the model computes three versions of its representation: a \emph{query}, a \emph{key}, and a \emph{value}.
  \item For each head, it measures how similar the current token's query is to every other token's key (via dot product).
  \item These similarities are turned into attention weights using softmax.
  \item The head output is a weighted sum of all value vectors, using these attention weights.
  \item All heads are concatenated and linearly mixed to form the final attention output.
\end{enumerate}

The bottleneck: storing and streaming all keys and values for every head becomes very expensive for long sequences and large models.

\section{Multi-Head Latent Attention (MLA)}

MLA keeps the \emph{queries} per head but \highlight{compresses \emph{keys} and \emph{values} into a shared low-dimensional latent space}, then reconstructs per-head K and V from that space. This greatly reduces KV-cache memory and bandwidth in autoregressive decoding and help the model process \highlight{long sequences} more efficiently.

\subsection{Latent K/V Definitions}

Let $d_{\ell} \ll d_{\text{model}}$ be the latent dimension. Starting from
\[
X \in \mathbb{R}^{L \times d_{\text{model}}},
\]
we compute:

\paragraph{Per-head queries.}
For each head $h$:
\[
Q^{(h)} = X W_Q^{(h)}, \quad W_Q^{(h)} \in \mathbb{R}^{d_{\text{model}} \times d_h}.
\]

\paragraph{Shared latent keys/values.}
We define head-independent latent K and V:
\begin{align*}
K_{\text{lat}} &= X W_{K,\text{lat}} \in \mathbb{R}^{L \times d_{\ell}}, \\
V_{\text{lat}} &= X W_{V,\text{lat}} \in \mathbb{R}^{L \times d_{\ell}},
\end{align*}
where $W_{K,\text{lat}}, W_{V,\text{lat}} \in \mathbb{R}^{d_{\text{model}} \times d_{\ell}}$.

\paragraph{Per-head reconstruction from latent space.}
For each head $h$, we have small ``decoder'' matrices:
\[
U_K^{(h)}, U_V^{(h)} \in \mathbb{R}^{d_{\ell} \times d_h},
\]
and reconstruct
\begin{align*}
K^{(h)} &= K_{\text{lat}} U_K^{(h)} \in \mathbb{R}^{L \times d_h}, \\
V^{(h)} &= V_{\text{lat}} U_V^{(h)} \in \mathbb{R}^{L \times d_h}.
\end{align*}

\paragraph{Attention computation.}
Attention itself is standard:
\[
\text{Attn}^{(h)}(X)
=
\text{softmax}\!\left(
\frac{Q^{(h)} (K^{(h)})^{\top}}{\sqrt{d_h}}
\right) V^{(h)}.
\]

In \highlight{autoregressive decoding, only $K_{\text{lat}}$ and $V_{\text{lat}}$ are cached}. The full per-head $K^{(h)}, V^{(h)}$ are reconstructed on the fly via $U_K^{(h)}, U_V^{(h)}$.

\subsection{MLA with RoPE-Aware and NoPE Components (High-Level)}

DeepSeek further splits the representation into:
\begin{itemize}
  \item A RoPE-aware part $(q_{\text{rope}}^{(h)}, k_{\text{rope}})$, which receives rotary position embeddings.
  \item A NoPE (no positional embedding) part $(q_{\text{nope}}^{(h)}, k_{\text{nope}}^{(h)})$, which stays purely semantic.
\end{itemize}

The attention score for head $h$ is then based on a combination of these components, typically written schematically as
\[
\alpha^{(h)}_{ij} \propto
\left(
\langle \tilde{q}^{(h)}_{\text{rope}, i}, \tilde{k}_{\text{rope}, j} \rangle
+
\langle q^{(h)}_{\text{nope}, i}, k^{(h)}_{\text{nope}, j} \rangle
\right),
\]
where $\tilde{q}$ and $\tilde{k}$ denote RoPE-rotated vectors (see Section~\ref{sec:rope}).

\subsection{Plain-Language Explanation of MLA}

\begin{enumerate}[label=\arabic*)]
  \item Instead of building a big key and value tensor for each head, the model builds a single compact \emph{latent} key tensor $K_{\text{lat}}$ and a single compact \emph{latent} value tensor $V_{\text{lat}}$ that all heads share.
  \item During generation, only $K_{\text{lat}}$ and $V_{\text{lat}}$ are stored in the KV-cache; this dramatically shrinks the memory footprint.
  \item Each attention head has its own small decoder matrices that convert the latent representation into its head-specific keys and values.
  \item The attention mechanism (dot product, softmax, weighted sum of values) stays unchanged; the only difference is that keys and values are factored through a low-rank latent space.
  \item As a result, DeepSeek can handle much longer sequences and larger models while keeping inference efficient.
\end{enumerate}

\section{Rotary Positional Embeddings (RoPE)}
\label{sec:rope}

RoPE injects positional information into queries and keys by applying position-dependent 2D rotations in feature space.

\subsection{Mathematical Definition}

Consider a per-head vector $x \in \mathbb{R}^{d_h}$ (either a query or a key). Group its coordinates into 2D pairs:
\[
x =
(x_0, x_1, x_2, x_3, \dots, x_{d_h-2}, x_{d_h-1}),
\]
with pairs
\[
(x_{2j}, x_{2j+1}), \quad j = 0, 1, \dots, d_h/2 - 1.
\]

For token position $i$ (0-indexed), define a frequency $\omega_j$ and an angle
\[
\omega_j = 10000^{-2j/d_h}, \qquad
\theta_{i,j} = i \cdot \omega_j.
\]

The RoPE-transformed pair at position $i$ is
\[
\begin{bmatrix}
\tilde{x}_{2j} \\
\tilde{x}_{2j+1}
\end{bmatrix}
=
\begin{bmatrix}
\cos \theta_{i,j} & -\sin \theta_{i,j} \\
\sin \theta_{i,j} & \cos \theta_{i,j}
\end{bmatrix}
\begin{bmatrix}
x_{2j} \\
x_{2j+1}
\end{bmatrix}.
\]

We denote this operation by
\[
\tilde{x} = \text{RoPE}(x, i).
\]

Queries and keys are then replaced by their RoPE-rotated versions:
\[
\tilde{q}_i^{(h)} = \text{RoPE}(q_i^{(h)}, i), \qquad
\tilde{k}_i^{(h)} = \text{RoPE}(k_i^{(h)}, i).
\]

The attention scores use $\tilde{q}$ and $\tilde{k}$:
\[
\text{score}^{(h)}_{ij}
=
\frac{
\tilde{q}_i^{(h)} \cdot \tilde{k}_j^{(h)}
}{
\sqrt{d_h}
}.
\]

A key property of RoPE is that $\tilde{q}_i^{\top}\tilde{k}_j$ depends on the \emph{relative} position $(i - j)$, enabling good extrapolation to long contexts.

\subsection{Plain-Language Explanation of RoPE}

\begin{enumerate}[label=\arabic*)]
  \item View each pair of adjacent dimensions of a query or key vector as a point in a 2D plane. So, \highlight{traditional sinusoidal PE is additive; ROPE is rotational}. 
  \item For a token at position $i$, rotate every such 2D point by a position-dependent angle; different pairs use different rotation rates (frequencies).
  \item When two rotated vectors from positions $i$ and $j$ are compared via dot product, the result now encodes both content similarity \emph{and} the relative position $(i-j)$.
  \item This gives the model a smooth, periodic representation of position that generalizes to sequences longer than those seen in training.
  \item In DeepSeek's MLA, only part of the representation receives RoPE (RoPE part), while another part remains unrotated (NoPE), letting the model disentangle positional information from purely semantic dimensions.
\end{enumerate}

\section{Mixture-of-Experts (DeepSeekMoE) and Router}

In a standard Transformer, the post-attention feed-forward network (FFN) is a single dense MLP. DeepSeek replaces this with a \emph{sparse} Mixture-of-Experts (MoE): many experts, but only a small subset are active per token.

\subsection{Mathematical Formulation}

Suppose the hidden state entering the MoE at token $t$ is $u_t \in \mathbb{R}^{d_{\text{model}}}$. There are $N$ experts, each being an FFN:
\[
\text{FFN}_i: \mathbb{R}^{d_{\text{model}}} \to \mathbb{R}^{d_{\text{model}}}, \quad i = 1,\dots,N.
\]

\paragraph{Routing scores.}
The router computes a score for each expert:
\[
s_{i,t} = w_i^\top u_t + b_i, \quad i = 1,\dots,N,
\]
where \highlight{$w_i \in \mathbb{R}^{d_{\text{model}}}$} and $b_i \in \mathbb{R}$ are learnable parameters.

\paragraph{Softmax probabilities.}
\[
p_{i,t} =
\frac{\exp(s_{i,t})}{\sum_{j=1}^N \exp(s_{j,t})}.
\]

\paragraph{Top-$K$ gating.}
Let $\mathcal{T}_t$ be the set of indices of the  \highlight{top-$K$ experts for token $t$ } (based on $p_{i,t}$). Define the \highlight{ normalized gates}:
\[
g_{i,t} =
\begin{cases}
\dfrac{p_{i,t}}{\sum_{j \in \mathcal{T}_t} p_{j,t}}, & i \in \mathcal{T}_t, \\[6pt]
0, & \text{otherwise}.
\end{cases}
\]

\paragraph{Expert outputs and aggregation.}
Each selected expert produces an output:
\[
y_{i,t} = \text{FFN}_i(u_t).
\]
The MoE output is the \highlight{gated sum}:
\[
y_t = \sum_{i=1}^N g_{i,t} \, y_{i,t}
    = \sum_{i \in \mathcal{T}_t} g_{i,t} \, \text{FFN}_i(u_t).
\]

\paragraph{DeepSeekMoE refinements (conceptual).}
DeepSeek introduces:
\begin{itemize}
  \item \emph{Fine-grained experts:} experts may be internally segmented into smaller sub-experts to increase diversity.
  \item  \highlight{ \emph{Shared experts:} } a small subset of experts are always active (generalists) to capture common patterns and stabilize training.
  \item \emph{Improved load balancing:} routing and training procedures reduce expert collapse and balance tokens across experts without heavy auxiliary losses.
\end{itemize}

\subsection{Plain-Language Explanation of MoE and Router}

\begin{enumerate}[label=\arabic*)]
  \item Think of many small specialist networks (experts), each good at processing certain kinds of tokens or patterns.
  \item For each token, the router looks at its hidden state and scores how suitable each expert is.
  \item Only a few top-scoring experts (e.g.\ $K=2$ or $K=4$) are chosen for that token; all others are skipped.
  \item The token is passed through those selected experts, producing several candidate transformed vectors.
  \item The router then takes a weighted average of these candidate vectors to form the final MoE output for that token.
  \item Because each token uses only a small subset of experts, the model can have a very large total number of parameters (high capacity) without increasing the per-token compute cost too much.
\end{enumerate}

\section{Reinforcement Learning and Policy Innovation in DeepSeek}

Beyond architecture, DeepSeek models (especially DeepSeekMath and DeepSeek-R1) are trained with a distinctive reinforcement learning (RL) framework designed to improve reasoning while keeping training efficient. The key policy innovation is \emph{Group Relative Policy Optimization} (GRPO), a critic-free variant of PPO that uses \emph{relative} rewards within groups of samples.

\subsection{RL Training Pipeline for DeepSeek-R1}

Conceptually, DeepSeek-R1 is derived from a strong base model (e.g.\ DeepSeek-V3-Base) using several RL-related stages:

\begin{enumerate}[label=\arabic*)]
  \item \textbf{DeepSeek-R1-Zero (pure RL on the base model).}
  \begin{itemize}
    \item Start from the base model without any supervised fine-tuning.
    \item Use GRPO with rule-based rewards that check final answer correctness and output format (for math, coding, logic tasks).
    \item The model is trained purely via RL to solve verifiable problems; sophisticated reasoning behaviors such as longer thinking, self-reflection, and re-checking emerge automatically.
  \end{itemize}

  \item \textbf{DeepSeek-R1 with cold-start data.}
  \begin{itemize}
    \item Collect a small but high-quality dataset of long chain-of-thought (CoT) solutions with human-designed formatting.
    \item Supervised fine-tune the base model on this CoT data to obtain an initial actor that is more readable and user-friendly.
  \end{itemize}

  \item \textbf{Reasoning-oriented RL.}
  \begin{itemize}
    \item Apply large-scale GRPO-based RL on reasoning-intensive tasks (math, coding, science, logic).
    \item Rewards combine final-answer correctness with auxiliary terms such as language consistency, encouraging both accuracy and readable reasoning.
  \end{itemize}

  \item \textbf{Rejection-sampling SFT.}
  \begin{itemize}
    \item From the RL-trained checkpoint, generate multiple candidate trajectories per prompt.
    \item Keep only correct and readable trajectories (via rule-based or model-based judgments).
    \item Use these high-quality trajectories to build a larger supervised dataset and fine-tune again, improving general skills (writing, dialogue, etc.).
  \end{itemize}

  \item \textbf{RL for all scenarios.}
  \begin{itemize}
    \item Run a second RL stage that mixes reasoning tasks with more general tasks.
    \item Combine rule-based rewards (for verifiable problems) with neural preference models that score helpfulness and harmlessness.
    \item Rewards focus mainly on the final answer for helpfulness, while safety rewards consider the whole reasoning trace and summary.
  \end{itemize}
\end{enumerate}

Overall, DeepSeek-R1 cycles between RL and supervised fine-tuning, gradually refining both reasoning ability and user-facing quality.

\subsection{ \highlight{ Group Relative Policy Optimization (GRPO)}: Objective and Advantages}

GRPO is a policy gradient method related to PPO but designed for large LLMs:

\begin{itemize}
  \item It \textbf{removes the critic network}. Instead of learning a separate value function, GRPO estimates advantages from \emph{relative} scores within a group of responses.
  \item It \textbf{keeps PPO-style clipping} and a KL-penalty to stabilize updates.
  \item It is \textbf{memory-efficient} because no large critic model is required.
\end{itemize}

\subsubsection*{GRPO Objective}

Let $P(Q)$ be a distribution over questions $q$. For each question $q$, GRPO samples a \emph{group} of $G$ responses from an old policy $\pi_{\theta_{\text{old}}}$:
\[
\{ o_1, o_2, \dots, o_G \} \sim \pi_{\theta_{\text{old}}}(\cdot \mid q).
\]

Each response $o_i$ receives a \highlight{ scalar reward $r_i$ } (for example, $r_i = 1$ if the final answer is correct, $0$ otherwise, plus optional format/language bonuses). From these rewards, GRPO constructs a \emph{group-normalized advantage}:
\begin{align*}
\bar{r} &= \frac{1}{G} \sum_{j=1}^G r_j, \\
s_r &= \sqrt{ \frac{1}{G} \sum_{j=1}^G (r_j - \bar{r})^2 + \epsilon }, \\
A_i &= \frac{r_i - \bar{r}}{s_r},
\end{align*}
where $\epsilon$ is a small constant to avoid division by zero.

Define the \highlight{ probability ratio } between the new policy $\pi_\theta$ and the old policy $\pi_{\theta_{\text{old}}}$:
\[
\rho_i(\theta) = \frac{\pi_\theta(o_i \mid q)}{\pi_{\theta_{\text{old}}}(o_i \mid q)}.
\]

GRPO maximizes the \highlight{ clipped, group-averaged objective with a KL regularizer}:
\[
J_{\text{GRPO}}(\theta)
=
\mathbb{E}_{q \sim P(Q)}
\left[
\frac{1}{G}
\sum_{i=1}^G
\left(
\min\!\big(
\rho_i(\theta) A_i,\,
\operatorname{clip}(\rho_i(\theta), 1-\varepsilon, 1+\varepsilon)\,A_i
\big)
- \beta \, \text{KL}\big(\pi_\theta(\cdot\mid q)\,\|\,\pi_{\text{ref}}(\cdot\mid q)\big)
\right)
\right],
\]
where $\varepsilon$ and $\beta$ are hyperparameters and $\pi_{\text{ref}}$ is a reference policy (e.g.\ the base model or an SFT checkpoint).

Intuitively:
\begin{itemize}
  \item If a response performs better than the group average ($A_i > 0$), GRPO pushes up its probability.
  \item If it performs worse ($A_i < 0$), GRPO pushes its probability down.
  \item The clipping on $\rho_i(\theta)$ prevents overly large policy updates, in the spirit of PPO.
  \item The KL term keeps the policy close to a reference model, stabilizing training.
\end{itemize}

\subsubsection*{Plain-Language Explanation of GRPO}

\begin{enumerate}[label=\arabic*)]
  \item For each question, sample several candidate answers from the \emph{old} model.
  \item Score every answer with a scalar reward (e.g.\ correct vs.\ incorrect, plus formatting/language bonuses).
  \item Compare answers only \emph{within the same group}:
  \begin{itemize}
    \item Answers better than the group average get positive advantage.
    \item Answers worse than the group average get negative advantage.
  \end{itemize}
  \item Compute how much more or less likely the \highlight{new} model wants to make each answer relative to the old one (this is the probability ratio $\rho_i$).
  \item Increase the probability of above-average answers and decrease the probability of below-average answers, but clip the update so it does not move too far in one step.
  \item Add a penalty if the new model drifts too far from a reference model (via the KL term).
\end{enumerate}

This design replaces the complex value-function critic with a simple, group-based baseline and works well when \highlight{ rewards are sparse but verifiable (e.g.\ final answer correctness). }

\subsection{Reward Design and Emergent Reasoning}

DeepSeek's RL pipeline uses relatively simple yet carefully chosen reward signals:

\begin{itemize}
  \item \textbf{Accuracy rewards} that check if the final answer matches the ground truth:
  \begin{itemize}
    \item For math, require the answer in a specific format (e.g.\ in a box or after a tag) and compare exactly.
    \item For coding, compile and run tests; reward depends on passing test cases.
  \end{itemize}
  \item \textbf{Format rewards} that enforce \highlight{ structural patterns} such as:
  \[
  \langle\text{think}\rangle \text{ reasoning process } \langle/\text{think}\rangle
  \langle\text{answer}\rangle \text{ final answer } \langle/\text{answer}\rangle.
  \]
  \item \textbf{Language and style rewards} that encourage output in a target language with readable formatting (e.g.\ markdown, clear summaries).
  \item \textbf{Preference rewards} in later stages, where a learned reward model judges helpfulness and harmlessness on open-ended tasks.
\end{itemize}

A notable empirical outcome is  \highlight{complex reasoning behaviors emerge without being hard-coded}:

\begin{itemize}
  \item The model learns to spend more ``thinking time'' (longer reasoning traces) on harder problems.
  \item It develops self-reflection patterns, such as revisiting previous steps when it suspects an error.
  \item It explores alternative solution strategies, guided only by correctness and simple format constraints.
\end{itemize}

In this sense, DeepSeek's RL approach is a \highlight{policy innovation}: rather than prescribing a specific style of reasoning, the training only rewards \highlight{outcomes and minimal structure}, and lets the policy autonomously discover effective internal strategies. (Instructor's aside: outcome-focused training also seems to work in children's development.)

\subsection{Policy Innovation in a Nutshell}

From a higher-level viewpoint, DeepSeek's RL and policy design innovate along three axes:

\begin{enumerate}[label=\arabic*)]
  \item \textbf{Critic-free, group-relative optimization.} GRPO replaces traditional value-function critics with group-based baselines, making large-scale RL on LLMs more memory-efficient and easier to implement.
  \item \textbf{Outcome-focused reward for reasoning.} Rewards mainly depend on verifiable final answers and light formatting constraints, leaving the reasoning process flexible. This encourages the model to invent its own reasoning heuristics.
  \item \textbf{Multi-stage RL + SFT pipeline.} Alternating between RL (to explore reasoning strategies) and SFT (to distill good behaviors into cleaner policies) yields models that are both strong reasoners and user-friendly assistants.
\end{enumerate}

\section{Putting a DeepSeek Block Together}

For a single layer $\ell$ and token $t$ at position $i$, the forward flow is:

\begin{enumerate}[label=\arabic*)]
  \item \textbf{Input:} hidden state $h_t^{(\ell)}$.

  \item \textbf{LayerNorm + MLA projections:}
  \[
  \hat{h}_t = \text{LN}\!\big(h_t^{(\ell)}\big).
  \]
  From $\hat{h}_t$ compute:
  \begin{itemize}
    \item per-head queries $q_t^{(h)}$,
    \item contributions to latent keys/values $K_{\text{lat}}, V_{\text{lat}}$.
  \end{itemize}

  \item \textbf{RoPE application:}
  \begin{itemize}
    \item For the RoPE-aware subspace, apply $\text{RoPE}(\cdot, i)$ to queries and keys.
    \item The NoPE subspace remains unrotated, preserving pure semantic content.
  \end{itemize}

  \item \textbf{Attention:}
  \begin{itemize}
    \item For each head, reconstruct $K^{(h)}, V^{(h)}$ from $K_{\text{lat}}, V_{\text{lat}}$.
    \item Compute attention scores and outputs per head.
    \item Concatenate head outputs and project to get $a_t$.
  \end{itemize}

  \item \textbf{First residual connection:}
  \[
  h_t^{(\ell,\text{attn})} = h_t^{(\ell)} + a_t.
  \]

  \item \textbf{LayerNorm + MoE:}
  \[
  \hat{u}_t = \text{LN}\!\big(h_t^{(\ell,\text{attn})}\big).
  \]
  The router scores experts on $\hat{u}_t$, picks top-$K$, and aggregates their FFN outputs into $m_t$.

  \item \textbf{Second residual connection:}
  \[
  h_t^{(\ell+1)} = h_t^{(\ell,\text{attn})} + m_t.
  \]

  \item \textbf{Repeat for next layer.}
\end{enumerate}

Over many layers, MLA+RoPE handle efficient context integration, while the MoE provides high-capacity, specialized non-linear transformations; the RL training pipeline reshapes the overall policy to prioritize correct, readable, and aligned reasoning.

\section{Summary}

References: DeepSeek-AI, ``DeepSeek-V2: A Strong, Economical, and Efficient Mixture-of-Experts Language Model,'' 2024, arXiv:\allowbreak 2405.04434; DeepSeek-AI, ``DeepSeek-V3 Technical Report,'' 2024, arXiv:\allowbreak 2412.19437; DeepSeek-AI, ``DeepSeek-R1: Incentivizing Reasoning Capability in LLMs via Reinforcement Learning,'' 2025, arXiv:\allowbreak 2501.12948; Z. Shao et al., ``DeepSeekMath,'' 2024, arXiv:\allowbreak 2402.03300 (the GRPO paper); J. Su et al., ``RoFormer,'' 2021, arXiv:\allowbreak 2104.09864.


\begin{itemize}
  \item \textbf{Multi-Head Latent Attention (MLA)} factorizes keys and values through a low-dimensional shared latent space. Only latent K/V are cached, and per-head K/V are reconstructed on demand, greatly reducing memory and bandwidth.
  \item \textbf{RoPE} injects positional information by rotating 2D feature pairs in queries and keys with position-dependent angles. This yields a relative-position-aware similarity and supports long-context extrapolation.
  \item \textbf{DeepSeekMoE} replaces dense FFNs with a sparse mixture-of-experts: a router chooses a few experts per token, enabling very high parameter capacity with limited per-token computation.
  \item \textbf{GRPO-based RL} trains DeepSeek's policies using group-relative, critic-free policy optimization and simple, verifiable rewards. This encourages the emergence of strong reasoning behavior with efficient compute.
  \item A \textbf{DeepSeek block} combines MLA+RoPE for attention and DeepSeekMoE for the FFN, wrapped in LayerNorm and residual connections, while the RL pipeline (GRPO + staged SFT) turns the base architecture into a reasoning-focused model like DeepSeek-R1.
\end{itemize}

\end{document}
